\documentclass[a4paper,12pt]{article}

% Pacotes essenciais
\usepackage[utf8]{inputenc}
\usepackage[T1]{fontenc}
\usepackage[brazil]{babel}
\usepackage{geometry}
\usepackage{amsmath}
\usepackage{graphicx}
\usepackage{booktabs}
\usepackage{array}
\usepackage{enumitem}
\usepackage{titlesec}
\usepackage{fancyhdr}
\usepackage{hyperref}
\usepackage{xcolor}
\usepackage{float}

% Configuração de margens
\geometry{
    top=2.5cm,
    bottom=2.5cm,
    left=3cm,
    right=2.5cm,
    headheight=15pt
}

% Configuração de cabeçalho e rodapé
\pagestyle{fancy}
\fancyhf{}
\fancyhead[R]{\thepage}
\renewcommand{\headrulewidth}{0pt}

% Configuração de títulos
\titleformat{\section}{\large\bfseries}{\thesection}{1em}{}
\titleformat{\subsection}{\normalsize\bfseries}{\thesubsection}{1em}{}
\titleformat{\subsubsection}{\normalsize\bfseries\itshape}{\thesubsubsection}{1em}{}

% Configuração de links
\hypersetup{
    colorlinks=true,
    linkcolor=black,
    citecolor=black,
    urlcolor=blue
}

% Título
\title{\textbf{SafeGuard}\\[0.5em]
\large Sistema de Monitoramento e Detecção de Quedas para Idosos}
\author{
    Felipe Marques \and Gabriel Ribeiro \and Pedro Dias
}
\date{\today}

\begin{document}

\maketitle

\begin{center}
\textbf{Documento de Projeto Técnico}
\end{center}

\vspace{1cm}

\tableofcontents

\newpage

%==============================================================================
\section{Introdução e Contexto}
%==============================================================================

Com o envelecimento da população, a segurança de idosos em ambientes domésticos ou institucionais torna-se uma preocupação crítica. Quedas são uma das principais causas de lesões e complicações de saúde nessa faixa etária. Muitas vezes, a ausência de um socorro rápido pode agravar as consequências de uma queda.

Este projeto propõe o desenvolvimento de um protótipo IoT (\textit{Internet of Things}) acessível e escalável para detectar quedas em tempo real e notificar cuidadores ou familiares.

%==============================================================================
\section{Definição do Problema}
%==============================================================================

A principal questão que este projeto visa resolver é a \textbf{incerteza e a demora no socorro} após um idoso sofrer uma queda, especialmente quando ele está sozinho ou incapaz de acionar ajuda manualmente.

\subsection{Problemas Específicos}

\begin{itemize}[leftmargin=2cm]
    \item \textbf{Isolamento:} Idosos que moram sozinhos ficam vulneráveis em situações de emergência.
    \item \textbf{Latência de Comunicação:} A falta de um sistema automático faz com que o tempo entre o incidente e o socorro seja indeterminado, aumentando riscos de saúde.
    \item \textbf{Falta de Dados:} Não há registro histórico da frequência ou intensidade dos movimentos/queda para análise médica preventiva.
\end{itemize}

%==============================================================================
\section{Solução Proposta}
%==============================================================================

O \textbf{SafeGuard} consiste em uma pulseira inteligente (\textit{wearable}) equipada com sensores de movimento e velocidade. O dispositivo monitora os padrões de aceleração do usuário continuamente. Ao detectar um padrão brusco de movimento seguido de ausência de movimento (característico de uma queda), o sistema dispara um alerta automaticamente.

\subsection{Fluxo da Solução}

\begin{enumerate}[leftmargin=2cm]
    \item \textbf{Detecção:} O sensor na pulseira identifica variação brusca de velocidade/aceleração (queda).
    \item \textbf{Processamento:} O microcontrolador analisa os dados e confirma o evento.
    \item \textbf{Comunicação:} O dado é enviado via protocolo MQTT.
    \item \textbf{Visualização:} Os dados são armazenados no \textbf{InfluxDB} e visualizados em tempo real no \textbf{Grafana}.
    \item \textbf{Ação:} O \textbf{Node-RED} recebe o alerta e dispara uma notificação (simulada ou real) para o cuidador.
\end{enumerate}

%==============================================================================
\section{Stack Tecnológica}
%==============================================================================

Para a arquitetura do sistema, utiliza-se um \textit{stack} moderno, leve e amplamente utilizado no mercado de IoT:

\begin{itemize}[leftmargin=2cm]
    \item \textbf{MQTT} (\textit{Message Queuing Telemetry Transport}): Protocolo de comunicação leve, ideal para dispositivos IoT com largura de banda limitada. Será o barramento de mensagens central.

    \item \textbf{InfluxDB:} Banco de dados de séries temporais (\textit{Time Series Database}). Perfeito para armazenar leituras de sensores contínuos (velocidade, aceleração X, Y, Z) com alta performance.

    \item \textbf{Grafana:} Plataforma de análise e visualização interativa. Permite a criação de \textit{dashboards} para monitorar a saúde do dispositivo e histórico de alertas.

    \item \textbf{Node-RED:} Ferramenta de desenvolvimento baseada em fluxo para conectar dispositivos, APIs e serviços online. Atuará como o ``cérebro'' lógico para roteamento de alertas.
\end{itemize}

%==============================================================================
\section{Especificação de Hardware}
%==============================================================================

Como o projeto será desenvolvido em ambiente de simulação, o hardware foi selecionado com base nos componentes disponíveis no \textbf{Wokwi}.

\subsection{Microcontrolador}

\textbf{ESP32:} Escolhido por possuir Wi-Fi e Bluetooth integrados (necessários para o MQTT) e por ser o microcontrolador mais robusto disponível no Wokwi para este fim.

\subsection{Sensores (Simulação de Comportamento)}

No ambiente Wokwi, utiliza-se a API de código para simular o comportamento físico de um sensor de velocidade/aceleração.

\subsubsection{Lógica de Simulação}

Como o Wokwi possui limitações de sensores físicos 3D, utilizou-se o código C++ no ESP32 para gerar dados de \textbf{Acelerômetro (IMU -- Unidade de Medição Inercial)}.

\begin{itemize}[leftmargin=2cm]
    \item \textbf{Simulação de ``Velocidade'':} A velocidade será derivada matematicamente a partir da aceleração no código, ou simulam-se picos de aceleração para representar a queda.

    \item \textbf{Interação:} Utiliza-se um \textbf{Potenciômetro} ou um \textbf{Botão} no Wokwi para simular o ``evento de queda'' durante os testes.

    \begin{itemize}
        \item \textit{Exemplo:} Girar o potenciômetro para o máximo simula uma aceleração brusca (impacto), disparando o fluxo de alerta.
    \end{itemize}
\end{itemize}

%==============================================================================
\section{Arquitetura do Sistema}
%==============================================================================

O fluxo de dados segue a seguinte topologia:

\subsection{Publisher (ESP32 -- Wokwi)}

\begin{itemize}
    \item Lê o estado simulado do sensor/botão.
    \item Verifica condição de queda (ex: Aceleração $>$ Limite).
    \item Publica nos tópicos MQTT:
    \begin{itemize}
        \item \texttt{safeguard/queda/alerta} (Status: ON/OFF)
        \item \texttt{safeguard/sensores/velocidade} (Dados numéricos)
    \end{itemize}
\end{itemize}

\subsection{Broker MQTT (Mosquitto/HiveMQ -- Servidor)}

\begin{itemize}
    \item Recebe as mensagens e as distribui para os inscritos (\textit{Subscribers}).
\end{itemize}

\subsection{Subscriber \& Storage (Node-RED + InfluxDB)}

\begin{itemize}
    \item O Node-RED assina o tópico MQTT.
    \item Os dados numéricos são gravados no InfluxDB.
    \item Se o tópico de alerta for ``ON'', o Node-RED envia uma notificação (\textit{Debug node} ou E-mail simulado).
\end{itemize}

\subsection{Visualization (Grafana)}

\begin{itemize}
    \item Conecta ao InfluxDB.
    \item Exibe gráficos de velocidade ao longo do tempo.
    \item Exibe um painel de ``Status do Sistema'' (Normal vs. Queda Detectada).
\end{itemize}

%==============================================================================
\section{Detalhamento da Lógica de Detecção}
%==============================================================================

Para identificar a queda, o \textit{firmware} no ESP32 implementa uma lógica simplificada de detecção baseada em vetor:

\begin{enumerate}[leftmargin=2cm]
    \item \textbf{Leitura Contínua:} O sistema lê os valores de aceleração simulados (Eixos X, Y, Z).

    \item \textbf{Cálculo do Módulo:} Calcula-se a magnitude do vetor aceleração:
    \[
    A = \sqrt{x^2 + y^2 + z^2}
    \]

    \item \textbf{Comparação (\textit{Threshold}):}
    \begin{itemize}
        \item Se $A > \text{LIMITE\_ALTO}$ (ex: 3g): Indica impacto súbito (possível queda).
        \item Seguido de $A < \text{LIMITE\_BAIXO}$ (ex: 0,5g) por alguns segundos: Indica imobilidade pós-impacto.
    \end{itemize}

    \item \textbf{Ação:} Se as condições acima forem verdadeiras, publica \texttt{QUEDA\_DETECTADA} no MQTT.
\end{enumerate}

%==============================================================================
\section{Entregáveis do Projeto}
%==============================================================================

\begin{enumerate}[leftmargin=2cm]
    \item \textbf{Código-fonte (\textit{Firmware}):} Código \texttt{.ino} para ESP32 rodando no simulador Wokwi.
    \item \textbf{Fluxo Node-RED:} Exportação JSON do fluxo de automação.
    \item \textbf{Dashboard Grafana:} Configuração dos painéis de monitoramento.
    \item \textbf{Documentação Técnica:} Este documento e manual de uso da simulação.
\end{enumerate}

%==============================================================================
\section{Resumo das Tecnologias e Ferramentas}
%==============================================================================

\begin{table}[H]
\centering
\caption{Componentes do Sistema}
\label{tab:componentes}
\begin{tabular}{>{\bfseries}l l l}
\toprule
\textbf{Componente} & \textbf{Tecnologia/Ferramenta} & \textbf{Função no Projeto} \\
\midrule
Hardware & ESP32 (Wokwi) & Coleta de dados e envio via Wi-Fi \\
Sensor & Simulação Código + Potenciômetro & Simular aceleração/velocidade e evento de queda \\
Protocolo & MQTT & Transporte leve das mensagens entre hardware e cloud \\
Backend/Logic & Node-RED & Processamento de regras e disparo de alertas \\
Banco de Dados & InfluxDB & Armazenamento histórico das leituras (Série Temporal) \\
Frontend & Grafana & Visualização gráfica dos dados e status \\
\bottomrule
\end{tabular}
\end{table}

%==============================================================================
\section{Considerações Finais}
%==============================================================================

Este documento serve como guia inicial para o desenvolvimento do projeto SafeGuard. A próxima etapa consiste em configurar o ambiente Wokwi e subir as instâncias do Broker MQTT, Node-RED, InfluxDB e Grafana -- podendo estas serem instaladas localmente via Docker ou em nuvem.

O sistema proposto oferece uma solução acessível e escalável para um problema crítico de saúde pública, utilizando tecnologias consolidadas no mercado de IoT e permitindo futuras expansões como integração com serviços de emergência e análise preditiva de quedas.

\end{document}
